% AutoResearch competition technical proposal v2
% Reconstructed in the supplied science template workspace; XeLaTeX, A4 portrait.

\documentclass[10pt,a4paper,UTF8]{ctexart}

\usepackage[left=15mm,right=15mm,top=14mm,bottom=14mm,headheight=14pt,headsep=4mm,footskip=8mm]{geometry}
\usepackage{fontspec}
\setCJKmainfont[BoldFont=FandolHei-Regular]{FandolSong-Regular}
\setCJKsansfont{FandolHei-Regular}
\setCJKmonofont{FandolFang-Regular}
\setsansfont{TeX Gyre Heros}
\setmainfont{TeX Gyre Termes}
\setmonofont{Latin Modern Mono}

\usepackage{xcolor}
\definecolor{ARnavy}{HTML}{17324D}
\definecolor{ARteal}{HTML}{008C8C}
\definecolor{ARlight}{HTML}{F3F6F8}
\definecolor{ARmid}{HTML}{D7E1E8}
\definecolor{ARgray}{HTML}{5F6B76}
\definecolor{ARamber}{HTML}{C47A16}
\definecolor{ARred}{HTML}{A83A3A}
\definecolor{ARgreen}{HTML}{2F7D61}

\usepackage{tikz}
\usetikzlibrary{arrows.meta,positioning,shapes.geometric,fit,calc,backgrounds,matrix}
\usepackage{graphicx}
\usepackage{float}
\usepackage{tabularx}
\usepackage{booktabs}
\usepackage{array}
\usepackage{multirow}
\usepackage{enumitem}
\usepackage{pifont}
\usepackage[most]{tcolorbox}
\usepackage{fancyhdr}
\usepackage{amsmath}
\usepackage{url}
\usepackage{hyperref}
\usepackage{bookmark}

\hypersetup{
  colorlinks=true,
  linkcolor=ARnavy,
  citecolor=ARteal,
  urlcolor=ARteal,
  pdftitle={AutoResearch：证据约束的科研闭环基础设施},
  pdfauthor={AutoResearch 项目组},
  pdfsubject={XH-202619 技术方案},
  pdfkeywords={AI Scientist,Qwen,Science 125,科研状态契约,证据治理}
}

\ctexset{
  section={format=\Large\bfseries\sffamily\color{ARnavy},beforeskip=0pt,afterskip=4pt},
  subsection={format=\large\bfseries\sffamily\color{ARteal},beforeskip=4pt,afterskip=2pt},
  subsubsection={format=\normalsize\bfseries\sffamily\color{ARnavy},beforeskip=3pt,afterskip=1pt}
}
\setcounter{secnumdepth}{2}
\setcounter{tocdepth}{2}
\setlength{\parindent}{0pt}
\setlength{\parskip}{2.5pt}
\linespread{1.05}
\setlist[itemize]{leftmargin=1.4em,itemsep=1pt,topsep=2pt,parsep=0pt}
\setlist[enumerate]{leftmargin=1.7em,itemsep=1pt,topsep=2pt,parsep=0pt}
\renewcommand{\arraystretch}{1.14}
\emergencystretch=2em

\pagestyle{fancy}
\fancyhf{}
\fancyhead[L]{\footnotesize\sffamily\color{ARgray}\nouppercase{\leftmark}}
\fancyhead[R]{\footnotesize\sffamily\color{ARgray}XH-202619 · 赛道一 / 方向 1A}
\fancyfoot[L]{\footnotesize\sffamily\color{ARgray}AutoResearch · 证据约束的科研闭环基础设施}
\fancyfoot[R]{\footnotesize\sffamily\color{ARnavy}\thepage}
\renewcommand{\headrulewidth}{0.4pt}
\renewcommand{\footrulewidth}{0pt}
\renewcommand{\sectionmark}[1]{\markboth{\thesection\quad #1}{}}

\newtcolorbox{innovationbox}[1]{enhanced,breakable,colback=ARlight,colframe=ARteal,boxrule=0.8pt,arc=1mm,left=2mm,right=2mm,top=1.5mm,bottom=1.5mm,title={\sffamily\bfseries #1},coltitle=white,colbacktitle=ARteal,fonttitle=\small}
\newtcolorbox{evidencebox}[1]{enhanced,breakable,colback=white,colframe=ARnavy,boxrule=0.8pt,arc=1mm,left=2mm,right=2mm,top=1.5mm,bottom=1.5mm,title={\sffamily\bfseries #1},coltitle=white,colbacktitle=ARnavy,fonttitle=\small}
\newtcolorbox{boundarybox}[1]{enhanced,breakable,colback=ARlight,colframe=ARgray,boxrule=0.7pt,arc=1mm,left=2mm,right=2mm,top=1.5mm,bottom=1.5mm,title={\sffamily\bfseries #1},coltitle=white,colbacktitle=ARgray,fonttitle=\small}
\newtcolorbox{reviewbox}[1]{enhanced,breakable,colback=white,colframe=ARred,boxrule=0.9pt,arc=1mm,left=2mm,right=2mm,top=1.5mm,bottom=1.5mm,title={\sffamily\bfseries #1},coltitle=white,colbacktitle=ARred,fonttitle=\small}
\newtcolorbox{flatbox}{enhanced,colback=ARlight,colframe=ARmid,boxrule=0.5pt,arc=1mm,left=2mm,right=2mm,top=1.3mm,bottom=1.3mm}

\newcommand{\runproof}{\textcolor{ARgreen}{\bfseries 运行证据}}
\newcommand{\codeproof}{\textcolor{ARnavy}{\bfseries 代码契约}}
\newcommand{\pending}{\textcolor{ARamber}{\bfseries【待补实验】}}
\newcommand{\blocked}{\textcolor{ARred}{\bfseries 阻断}}
\newcommand{\cmark}{\textcolor{ARgreen}{\ding{51}}}
\newcommand{\unknown}{\textcolor{ARgray}{公开材料未明确}}
\newcommand{\metriccard}[2]{\begin{tcolorbox}[enhanced,colback=white,colframe=ARmid,boxrule=0.6pt,arc=1mm,left=1mm,right=1mm,top=1mm,bottom=1mm,halign=center]\sffamily\bfseries\Large\color{ARnavy}#1\\[-1mm]{\scriptsize\color{ARgray}#2}\end{tcolorbox}}
\newcommand{\pagetitle}[2]{\phantomsection\addcontentsline{toc}{section}{#1}\markboth{#1}{}\begin{tcolorbox}[enhanced,colback=ARnavy,colframe=ARnavy,boxrule=0pt,arc=0mm,left=3mm,right=3mm,top=2mm,bottom=2mm]\color{white}\sffamily\bfseries\Large #1\hfill{\small\mdseries #2}\end{tcolorbox}}

\tikzset{
  arnode/.style={rounded corners=1.5mm,draw=ARmid,fill=white,very thick,align=center,minimum height=8mm,font=\small\sffamily,text=ARnavy},
  arprimary/.style={arnode,draw=ARteal,fill=ARlight},
  arblock/.style={arnode,draw=ARred,fill=white,text=ARred},
  ardata/.style={-{Latex[length=2.2mm]},line width=1.0pt,draw=ARteal},
  arpermission/.style={-{Latex[length=2.2mm]},line width=0.9pt,draw=ARnavy,dashed},
  arfeedback/.style={-{Latex[length=2.2mm]},line width=0.9pt,draw=ARred,densely dotted},
  arlink/.style={-{Latex[length=2mm]},line width=0.8pt,draw=ARgray}
}

\begin{document}

% ==================== P1 COVER ====================
\pdfbookmark[0]{封面}{cover}
\thispagestyle{empty}
\begin{tikzpicture}[remember picture,overlay]
  \fill[ARnavy] (current page.north west) rectangle ([yshift=-54mm]current page.north east);
  \fill[ARteal] ([yshift=-54mm]current page.north west) rectangle ([yshift=-57mm]current page.north east);
  \draw[ARmid,line width=0.5pt] ([xshift=15mm,yshift=15mm]current page.south west) -- ([xshift=-15mm,yshift=15mm]current page.south east);
  \foreach \x in {0,...,8}{\draw[ARteal!25,line width=0.35pt] ([xshift=15mm+20mm*\x,yshift=38mm]current page.south west) -- ++(0,18mm);}
\end{tikzpicture}

\vspace*{9mm}
{\sffamily\color{white}\large 国家级人工智能科技竞赛 · 技术方案}

\vspace{24mm}
{\sffamily\bfseries\fontsize{30}{36}\selectfont\color{ARnavy}AutoResearch}

\vspace{4mm}
{\sffamily\bfseries\fontsize{21}{28}\selectfont\color{ARnavy}证据约束的科研闭环基础设施}

\vspace{5mm}
{\sffamily\fontsize{13}{20}\selectfont\color{ARteal}基于 Qwen、科研状态契约与三层门禁的\\Science 125 问 AI Scientist 技术方案}

\vspace{13mm}
\begin{tcolorbox}[enhanced,colback=ARlight,colframe=ARteal,boxrule=1pt,arc=1mm,left=4mm,right=4mm,top=4mm,bottom=4mm]
{\sffamily\bfseries\large\color{ARnavy}一句话价值主张}\par
\vspace{1mm}
把大模型的科研推理转化为\textbf{可追溯、可复核、可阻断、可迭代}的科研状态流，使“生成建议”升级为“依据真实证据推进下一步研究”。
\end{tcolorbox}

\vfill
\begin{tabularx}{\textwidth}{>{\sffamily\bfseries\color{ARnavy}}p{32mm}X}
项目编号 & XH-202619 \\
参赛方向 & 赛道一：科学问题 · 方向 1A 科学假设生成与研究计划设计 \\
项目团队 & AutoResearch 项目组 \\
版本日期 & 2026 年 8 月 \\
\end{tabularx}

\vspace{21mm}
{\small\color{ARgray}本方案严格区分“研究计划”与“真实观测”：125 份结构化计划不等于 125 项已执行实验。}
\clearpage

% ==================== P2 TOC ====================
\pdfbookmark[0]{目录}{toc}
\markboth{目录}{}
\begin{tcolorbox}[enhanced,colback=ARnavy,colframe=ARnavy,boxrule=0pt,arc=0mm,left=3mm,right=3mm,top=2mm,bottom=2mm]
\color{white}\sffamily\bfseries\Large 目录\hfill{\small\mdseries 评委阅读路径}
\end{tcolorbox}

\vspace{2mm}
\begin{minipage}[t]{0.64\textwidth}
\tableofcontents
\end{minipage}\hfill
\begin{minipage}[t]{0.32\textwidth}
\begin{evidencebox}{建议阅读路径}
\textbf{20 秒}\quad P3\par
定位、成果、边界\par\medskip
\textbf{1 分钟}\quad P4\par
六项核心创新\par\medskip
\textbf{5 分钟}\quad P5--P10\par
架构、状态、闭环、门禁\par\medskip
\textbf{10 分钟}\quad P11--P16\par
125 问、q001、评估与证据索引
\end{evidencebox}
\begin{boundarybox}{证据状态图例}
\runproof：已有运行产物\par
\codeproof：实现已核验，效果未必实测\par
\pending：尚无可报告结果\par
\blocked：门禁拒绝继续流转
\end{boundarybox}
\end{minipage}

\vfill
\begin{flatbox}
\textbf{文档导航说明：}目录条目、章节标题和参考文献均建立 PDF 超链接与书签。主文保留评委作出判断所需的机制和证据；完整代码路径、运行目录、指标与补证协议集中于附录。
\end{flatbox}
\clearpage

% ==================== P3 OVERVIEW ====================
\pagetitle{项目总览}{一图读懂}

\begin{innovationbox}{项目定位}
\textbf{AutoResearch 不是“自动写论文”的聊天机器人，而是一套证据优先的 AI Scientist 科研状态机：}以 Qwen 承担规模化推理，以临时智能体与 Skills 执行方法，以证据哈希、checkpoint、数值/引用守卫和三层门禁约束每一次状态跃迁。
\end{innovationbox}

\begin{minipage}[t]{0.32\textwidth}
\begin{flatbox}\textbf{问题 1｜证据漂移}\par 文献事实、模型推断与实验观测容易混写。\end{flatbox}
\end{minipage}\hfill
\begin{minipage}[t]{0.32\textwidth}
\begin{flatbox}\textbf{问题 2｜过程不可复核}\par 多轮生成覆盖中间状态与失败原因。\end{flatbox}
\end{minipage}\hfill
\begin{minipage}[t]{0.32\textwidth}
\begin{flatbox}\textbf{问题 3｜自动化无边界}\par 结构合规、科学合理与责任授权被混为一谈。\end{flatbox}
\end{minipage}

\vspace{2mm}
\begin{figure}[H]
\centering
\begin{tikzpicture}[node distance=3mm]
  \node[arprimary,minimum width=21mm] (q) {科学问题};
  \node[arnode,minimum width=22mm,right=of q] (e) {证据锁定};
  \node[arnode,minimum width=22mm,right=of e] (h) {候选假设};
  \node[arnode,minimum width=22mm,right=of h] (p) {可执行计划};
  \node[arnode,minimum width=25mm,right=of p] (x) {真实 pilot\\或待执行协议};
  \node[arnode,minimum width=22mm,right=of x] (r) {独立评审};
  \node[arblock,minimum width=22mm,right=of r] (d) {交付 / 阻断};
  \draw[ardata] (q)--(e); \draw[ardata] (e)--(h); \draw[ardata] (h)--(p);
  \draw[ardata] (p)--(x); \draw[ardata] (x)--(r); \draw[ardata] (r)--(d);
  \draw[arfeedback] (d.south) to[bend right=22] node[below,font=\scriptsize\sffamily,text=ARred]{评审问题回灌为下一轮任务} (e.south);
\end{tikzpicture}
\par\small\color{ARgray}最小可信科研闭环：数据流为实线，负面反馈为点线；阻断是正式结果，不是异常垃圾。\normalsize
\end{figure}

\begin{minipage}[t]{0.32\textwidth}
\begin{flatbox}\textbf{01 科研状态契约}\par 稳定 ID＋内容哈希绑定证据、计划、观测与决策。\end{flatbox}
\begin{flatbox}\textbf{04 不可变科研闭环}\par 十二阶段 checkpoint，失败另存 artifact。\end{flatbox}
\end{minipage}\hfill
\begin{minipage}[t]{0.32\textwidth}
\begin{flatbox}\textbf{02 临时智能体＋Skills}\par 按任务创建、最小授权、结构回执、执行后清理。\end{flatbox}
\begin{flatbox}\textbf{05 数值/引用＋三层门禁}\par 无法溯源即拒绝；结构、科学、责任分层判定。\end{flatbox}
\end{minipage}\hfill
\begin{minipage}[t]{0.32\textwidth}
\begin{flatbox}\textbf{03 原始证据/派生记忆分离}\par 摘要只导航，不能升级为科学证据。\end{flatbox}
\begin{flatbox}\textbf{06 规模化任务编译}\par 125 问转为带数据、基线、指标和终止条件的计划。\end{flatbox}
\end{minipage}

\vspace{2mm}
\begin{minipage}{0.135\textwidth}\metriccard{125/125}{结构化研究计划}\end{minipage}\hfill
\begin{minipage}{0.135\textwidth}\metriccard{125}{Markdown}\end{minipage}\hfill
\begin{minipage}{0.135\textwidth}\metriccard{125}{PDF}\end{minipage}\hfill
\begin{minipage}{0.135\textwidth}\metriccard{1}{q001 真实 pilot}\end{minipage}\hfill
\begin{minipage}{0.135\textwidth}\metriccard{4}{零模型}\end{minipage}\hfill
\begin{minipage}{0.135\textwidth}\metriccard{major}{revision}\end{minipage}\hfill
\begin{minipage}{0.135\textwidth}\metriccard{已阻断}{自动交付}\end{minipage}

\begin{boundarybox}{当前能力边界}
125 份是\textbf{研究计划}，不是 125 项已执行实验；q050、q091、q119 为领域定制计划态样例。q001 是探索性计算，\textbf{没有证明任何素数猜想}。Qwen 是当前核心，端点与供应商由配置替换。
\end{boundarybox}

\begin{minipage}[t]{0.49\textwidth}
\begin{evidencebox}{当前状态快照｜已经发生}
\begin{itemize}
  \item 125 题均形成结构化计划与双格式交付；
  \item q001 完成真实探索性计算并锁定原始产物；
  \item 独立评审发现 2 项 major、2 项 minor 问题；
  \item 阻断回执被保存，自动交付未发生。
\end{itemize}
\end{evidencebox}
\end{minipage}\hfill
\begin{minipage}[t]{0.49\textwidth}
\begin{boundarybox}{下一步验证｜尚未发生}
\begin{itemize}
  \item 其余 124 题需按数据与适配器条件逐题执行；
  \item 125 题内部 JSON Schema 批量复核待完成；
  \item 数字/引用对抗、恢复、成本与消融均待实验；
  \item Skill 影子评估已有契约，真实晋级链待补。
\end{itemize}
\end{boundarybox}
\end{minipage}

\begin{innovationbox}{20 秒判断句}
\textbf{输入不是一段提示词，输出也不只是一篇文本：}输入被编译为科研对象，输出携带证据绑定、状态、版本与门禁结论；系统既能推进，也能在证据不足时留下可解释的停止点。
\end{innovationbox}
\clearpage

% ==================== P4 CHALLENGES / GOALS / INNOVATIONS ====================
\section{赛题挑战与现有方案缺口}
\begin{tabularx}{\textwidth}{>{\bfseries\color{ARnavy}}p{28mm}X X}
\toprule
挑战 & 常见失效方式 & AutoResearch 的控制点 \\
\midrule
从开放问题到可证伪任务 & 输出泛化建议，缺变量、对照、判据和停止条件 & 任务编译 Schema 固化科学目标、竞争解释、数据、基线、指标与资源 \\
证据进入推理链 & 检索内容、模型记忆和真实观测混写 & 锁定引用目录；原始证据不可覆盖；派生记忆显式标记为非证据 \\
反馈真正改变状态 & 评审只是文末意见，下一轮生成可绕过失败 & checkpoint＋阻断回执；评审结论决定可交付、修订或停止 \\
\bottomrule
\end{tabularx}

\section{核心设计目标}
\begin{minipage}[t]{0.24\textwidth}\begin{flatbox}\textbf{G1 可检查}\par 每个结论可定位到对象、文件与哈希。\end{flatbox}\end{minipage}\hfill
\begin{minipage}[t]{0.24\textwidth}\begin{flatbox}\textbf{G2 可验证}\par 假设附带对照、指标和反证条件。\end{flatbox}\end{minipage}\hfill
\begin{minipage}[t]{0.24\textwidth}\begin{flatbox}\textbf{G3 可迭代}\par 评审与观测形成下一轮显式输入。\end{flatbox}\end{minipage}\hfill
\begin{minipage}[t]{0.24\textwidth}\begin{flatbox}\textbf{G4 可负责}\par 高风险执行和交付必须经过责任门禁。\end{flatbox}\end{minipage}

\section{核心创新总览}
\small
\begin{tabularx}{\textwidth}{>{\bfseries\color{ARnavy}}p{30mm}p{34mm}X p{27mm}}
\toprule
一级创新 & 解决的问题 & 核心机制与差异 & 当前证据 \\
\midrule
证据绑定的科研状态契约 & 对话与摘要不可稳定复核 & 状态对象＋稳定 ID＋哈希；计划/观测、原始/派生分别建模 & \runproof\,＋\codeproof \\
临时智能体与可审计 Skills & 常驻 Agent 权限与上下文污染 & 按任务装配权限、证据、Skill、Schema 和预算；保留路由理由与哈希 & \runproof\,＋\codeproof \\
受控 Skill 进化 & 错误经验直接写回系统 & 开发/隐藏证据分区、影子评估、晋级与回滚；安全规则不可自改 & \codeproof；\pending \\
十二阶段不可变闭环 & 失败被重生成覆盖 & 阶段 checkpoint、上游哈希引用、failure artifact 与反馈回灌 & \runproof\,＋\codeproof \\
数值/引用守卫与三层门禁 & 编造数字、假引用、责任不清 & 数值占位绑定、引用精确投影；机器/科学/人类分层判定 & 前两层有\runproof \\
确定性交付与任务编译 & 文本难执行、多格式漂移 & 同源渲染 MD/TeX/PDF/manifest；125 问编译为研究协议 & \runproof\,＋\codeproof \\
\bottomrule
\end{tabularx}
\normalsize

\begin{innovationbox}{为什么这不是普通 RAG 或多智能体编排}
RAG 决定“取回什么”，多智能体框架决定“谁与谁通信”；AutoResearch 进一步决定\textbf{什么可被认定为证据、当前处于何种科研状态、能否跃迁、失败如何保留、谁对执行和交付负责}。差异化来自状态与门禁契约，而非 Agent 数量或提示词长度。
\end{innovationbox}

\subsection*{创新证据状态：不把设计愿景写成运行结论}
\begin{minipage}[t]{0.32\textwidth}\begin{evidencebox}{运行级}
q001 证据状态、checkpoint、临时交互、真实 pilot、独立评审与阻断；125 题双格式交付。
\end{evidencebox}\end{minipage}\hfill
\begin{minipage}[t]{0.32\textwidth}\begin{evidencebox}{代码级}
状态 Schema、原始/派生分离、最小权限、数值/引用守卫、确定性渲染与 Skill 影子晋级契约。
\end{evidencebox}\end{minipage}\hfill
\begin{minipage}[t]{0.32\textwidth}\begin{boundarybox}{待验证级}
受控 Skill 晋级效果、批量守卫性能、checkpoint 恢复、跨模型消融、成本与人工修改量。
\end{boundarybox}\end{minipage}

\begin{boundarybox}{一级创新的判定规则}
本方案只把有代码或运行对象的机制列为核心创新；运行效果尚未测量时，明确保留【待补实验】，不以架构图替代证据。
\end{boundarybox}
\clearpage

% ==================== P5 ARCHITECTURE ====================
\section{系统总体架构}
\subsection{五层架构与三类流}
\begin{figure}[H]
\centering
\begin{tikzpicture}[x=1cm,y=1cm]
  \node[arprimary,minimum width=14.7cm,minimum height=1.25cm] (l1) at (0,5.7) {\textbf{交互层}\quad Science 125 问 · 数据/仪器入口 · 模型配置 · 人类反馈与授权};
  \node[arnode,minimum width=14.7cm,minimum height=1.35cm] (l2) at (0,3.9) {\textbf{编排层}\quad 科研状态机 · 十二阶段 checkpoint · 重试/恢复 · 机器门禁 · 科学门禁 · 责任门禁};
  \node[arnode,minimum width=14.7cm,minimum height=1.65cm] (l3) at (0,1.75) {\textbf{临时智能体层}\quad 目标收窄｜文献｜假设｜批评｜规划｜执行｜分析｜独立评审｜写作投影\\[-0.5mm]{\scriptsize 每个任务独立角色、证据子集、输出 Schema、上下文预算与清理回执}};
  \node[arnode,minimum width=14.7cm,minimum height=1.45cm] (l4) at (0,-0.55) {\textbf{Skills 层}\quad 认知 Skill：反证/创新性审查\quad 方法 Skill：统计/仿真/实验协议\quad 治理 Skill：引用/数字/权限};
  \node[arnode,minimum width=14.7cm,minimum height=1.7cm] (l5) at (0,-2.95) {\textbf{证据与基础设施层}\quad 原始文献/数据/日志/manifest · Obsidian 兼容知识库 · 版本库 · 沙箱\\[-0.5mm]Qwen 当前推理核心｜OpenAI-compatible 配置接口｜领域适配器};

  \node[arblock,minimum width=2.25cm,minimum height=1cm] (g1) at (6.1,3.9) {结构/溯源\\阻断点};
  \node[arblock,minimum width=2.25cm,minimum height=1cm] (g2) at (6.1,1.75) {独立评审\\阻断点};
  \node[arblock,minimum width=2.25cm,minimum height=1cm] (g3) at (6.1,-0.55) {权限/安全\\阻断点};

  \draw[ardata] (l1.south west) ++(1.0,0) -- ++(0,-0.55) -- ([xshift=1cm]l2.north west);
  \draw[ardata] (l2.south west) ++(1.0,0) -- ++(0,-0.65) -- ([xshift=1cm]l3.north west);
  \draw[ardata] (l3.south west) ++(1.0,0) -- ++(0,-0.5) -- ([xshift=1cm]l4.north west);
  \draw[ardata] (l4.south west) ++(1.0,0) -- ++(0,-0.65) -- ([xshift=1cm]l5.north west);
  \draw[arpermission] ([xshift=-1cm]l2.south east) -- ([xshift=-1cm]l3.north east);
  \draw[arpermission] ([xshift=-1cm]l3.south east) -- ([xshift=-1cm]l4.north east);
  \draw[arfeedback] ([xshift=0.2cm]l5.north) to[bend right=24] ([xshift=0.2cm]l2.south);

  \node[font=\scriptsize\sffamily,text=ARteal,anchor=west] at (-7.35,-4.35) {\rule{7mm}{1pt}\ 数据/产物流};
  \node[font=\scriptsize\sffamily,text=ARnavy,anchor=west] at (-3.5,-4.35) {-- -- -- 权限流};
  \node[font=\scriptsize\sffamily,text=ARred,anchor=west] at (0.1,-4.35) {······ 反馈/阻断流};
\end{tikzpicture}
\caption{AutoResearch 总体分层架构。数据、权限和反馈采用统一线型；红色节点是可真实失败的控制点。}
\label{fig:architecture-v2}
\end{figure}

\begin{minipage}[t]{0.32\textwidth}\begin{evidencebox}{Qwen 的位置}当前推理核心，承担问题收窄、假设生成、计划修订等认知任务；模型回执记录 provider 与 model。\end{evidencebox}\end{minipage}\hfill
\begin{minipage}[t]{0.32\textwidth}\begin{evidencebox}{可替换边界}代码从配置读取 \texttt{base\_url}、\texttt{api\_key}、\texttt{model\_name}；密钥不进入产物与版本库。\end{evidencebox}\end{minipage}\hfill
\begin{minipage}[t]{0.32\textwidth}\begin{boundarybox}{未宣称内容}尚无跨供应商同题质量、时延与成本对照；“可替换”是接口能力，不是性能等价结论。\end{boundarybox}\end{minipage}

\subsection*{跨层约束与可检查输出}
\small
\begin{tabularx}{\textwidth}{>{\bfseries\color{ARnavy}}p{26mm}p{43mm}X}
\toprule
约束面 & 在哪一层实施 & 评委可检查的输出 \\
\midrule
数据边界 & 证据层锁定文件、引用与哈希；编排层验证绑定 & manifest、verified files、引用目录状态、输入/产物哈希 \\
权限边界 & 编排层下发角色策略；Agent/Skill 层只获得必要工具 & 任务 Schema、Skill manifest、允许的证据子集、清理/紧急回执 \\
状态边界 & 编排层状态机区分计划、执行、观测、评审与交付 & checkpoint、parent hash、状态字段、failure/block receipt \\
责任边界 & 机器、科学与人类门禁各自独立作出判定 & 字段级拒绝、独立 recommendation、授权/保持/拒绝记录 \\
\bottomrule
\end{tabularx}
\normalsize

\begin{innovationbox}{架构判定原则}
任何智能体都不能仅凭自然语言“宣布成功”改变科研状态；状态提升必须由编排层读取结构化回执和已验证证据，并满足相应门禁条件。
\end{innovationbox}
\clearpage

% ==================== P6 EVIDENCE / CONTRACTS ====================
\section{证据图谱与科研状态契约}
\subsection{对象关系：让每个结论有来源、状态和版本}
\begin{figure}[H]
\centering
\begin{tikzpicture}[x=1cm,y=1cm]
  \node[arprimary,minimum width=2.6cm] (goal) at (-6.3,3.6) {科学目标\\{\ttfamily goal\_id}};
  \node[arnode,minimum width=2.7cm] (lit) at (-6.3,1.6) {文献证据\\DOI/URL＋hash};
  \node[arnode,minimum width=2.7cm] (data) at (-6.3,-0.4) {原始数据文件\\path＋SHA-256};
  \node[arnode,minimum width=2.7cm] (log) at (-6.3,-2.4) {执行日志\\run\_id＋hash};

  \node[arnode,minimum width=3.0cm] (hyp) at (-2.0,3.6) {候选假设\\可证伪判据};
  \node[arnode,minimum width=3.0cm] (plan) at (-2.0,1.2) {研究计划\\plan-only / executable};
  \node[arnode,minimum width=3.0cm] (obs) at (-2.0,-1.2) {观测对象\\metrics＋manifest};
  \node[arnode,minimum width=3.0cm] (derived) at (-2.0,-3.2) {派生记忆\\非科学证据};

  \node[arnode,minimum width=3.0cm] (review) at (2.5,1.2) {独立评审\\recommendation};
  \node[arblock,minimum width=3.0cm] (decision) at (2.5,-1.2) {决策对象\\交付 / 修订 / 阻断};
  \node[arnode,minimum width=3.0cm] (version) at (6.5,0) {版本对象\\parent hash＋diff};

  \draw[ardata] (goal)--(hyp); \draw[ardata] (lit)--(hyp); \draw[ardata] (hyp)--(plan);
  \draw[ardata] (data)--(obs); \draw[ardata] (log)--(obs); \draw[ardata] (plan)--(obs);
  \draw[ardata] (plan)--(review); \draw[ardata] (obs)--(review); \draw[ardata] (review)--(decision);
  \draw[ardata] (decision)--(version); \draw[arfeedback] (version.north) to[bend left=25] (hyp.east);
  \draw[arpermission] (derived)--(plan) node[midway,right,font=\scriptsize\sffamily,text=ARnavy]{仅导航};

  \node[draw=ARteal,dashed,rounded corners,fit=(lit)(data)(log),inner sep=4mm,label={[font=\scriptsize\sffamily,text=ARteal]above:不可覆盖的原始证据域}] {};
  \node[draw=ARnavy,dashed,rounded corners,fit=(hyp)(plan)(obs),inner sep=4mm,label={[font=\scriptsize\sffamily,text=ARnavy]above:科研状态域}] {};
\end{tikzpicture}
\caption{证据对象关系。实线表示数据/产物绑定，虚线表示受限使用，点线表示反馈生成新版本。}
\label{fig:evidence-objects}
\end{figure}

\begin{innovationbox}{科研状态契约}
一个可信状态可抽象为 $S_t=\langle id, type, payload, evidence\_bindings, parent\_hash, artifact\_hash, status\rangle$。下游只能引用已验证的 $S_t$；任何字段变化都会产生新哈希和新版本，不能静默覆盖原状态。
\end{innovationbox}

\small
\begin{tabularx}{\textwidth}{>{\bfseries\color{ARnavy}}p{30mm}X X}
\toprule
强制分离 & 机器约束 & 科学意义 \\
\midrule
原始证据 / 派生记忆 & \texttt{derived\_context\_is\_evidence=false}；原始字节保留、派生内容可重建 & 摘要可帮助检索，但不能替代文献、数据或运行记录 \\
计划态 / 观测态 & 无兼容适配器时写入“未执行”；交付校验拒绝虚构结果图和数字 & 125 份计划不会被误报为 125 次实验 \\
作者修订 / 独立评审 & 评审器不重写作者产物，另存 recommendation 与证据清单 & 避免同一模型链既当作者又替自己放行 \\
\bottomrule
\end{tabularx}
\normalsize

\begin{evidencebox}{已核验实例}
q001 独立评审 artifact 记录 18 个已验证文件及文件集合哈希；引用目录检查结果为 \texttt{verified\_exact\_union}（精确并集）。评审与阻断回执均作为新对象保存。
\end{evidencebox}
\clearpage

% ==================== P7 TEMP AGENTS / SKILLS ====================
\section{临时智能体与可审计 Skills}
\subsection{按任务创建、按角色授权、完成后销毁}
\begin{figure}[H]
\centering
\begin{tikzpicture}[x=1cm,y=1cm]
  \node[arprimary,minimum width=2.1cm] (a1) at (-7.2,2.7) {1 任务编译};
  \node[arnode,minimum width=2.3cm] (a2) at (-4.4,2.7) {2 角色授权\\最小权限};
  \node[arnode,minimum width=3.0cm] (a3) at (-1.0,2.7) {3 上下文装配\\固定顺序＋预算};
  \node[arnode,minimum width=2.5cm] (a4) at (2.6,2.7) {4 受限执行\\工具白名单};
  \node[arnode,minimum width=2.4cm] (a5) at (5.7,2.7) {5 Schema 回执\\输出哈希};
  \node[arblock,minimum width=2.2cm] (a6) at (7.1,0.2) {6 清理/销毁\\审计回执};
  \draw[ardata] (a1)--(a2); \draw[arpermission] (a2)--(a3); \draw[ardata] (a3)--(a4); \draw[ardata] (a4)--(a5); \draw[ardata] (a5)--(a6);

  \node[arnode,minimum width=2.25cm] (c1) at (-5.7,0.1) {角色政策};
  \node[arnode,minimum width=2.25cm] (c2) at (-3.2,0.1) {任务字节};
  \node[arnode,minimum width=2.25cm] (c3) at (-0.7,0.1) {允许 Skills};
  \node[arnode,minimum width=2.25cm] (c4) at (1.8,0.1) {锁定证据};
  \node[arnode,minimum width=2.25cm] (c5) at (4.3,0.1) {派生导航};
  \draw[arpermission] (c1.north)--(a3.south); \draw[arpermission] (c2.north)--(a3.south);
  \draw[arpermission] (c3.north)--(a3.south); \draw[arpermission] (c4.north)--(a3.south); \draw[arpermission] (c5.north)--(a3.south);

  \node[arprimary,minimum width=2.5cm] (s1) at (-6.0,-2.4) {文献＋pilot\\真实证据};
  \node[arnode,minimum width=2.5cm] (s2) at (-2.8,-2.4) {候选 Skill\\版本＋哈希};
  \node[arnode,minimum width=3.0cm] (s3) at (0.8,-2.4) {隐藏证据影子评估\\不接触测试标签};
  \node[arnode,minimum width=2.1cm] (s4) at (4.3,-2.4) {晋级};
  \node[arblock,minimum width=2.1cm] (s5) at (7.0,-2.4) {拒绝/回滚};
  \draw[ardata] (s1)--(s2); \draw[ardata] (s2)--(s3); \draw[ardata] (s3)--(s4); \draw[arfeedback] (s3)--(s5);
  \node[font=\scriptsize\sffamily,text=ARred] at (0.8,-3.45) {安全策略、审批门禁与发布规则不可由 Skill 自我修改};
\end{tikzpicture}
\caption{临时智能体与 Skill 生命周期。上：最小权限任务执行；下：受控 Skill 进化设计。}
\label{fig:agent-skill-lifecycle}
\end{figure}

\begin{tabularx}{\textwidth}{>{\bfseries\color{ARnavy}}p{27mm}X p{48mm}}
\toprule
Skill 类型 & 作用 & 审计字段 \\
\midrule
认知 Skill & 收窄问题、设计反证、创新性三角审查、竞争解释 & 路由理由、来源、版本、内容哈希 \\
方法 Skill & 统计检验、数据处理、仿真/实验协议与终止条件 & 输入/输出 Schema、工具权限、证据绑定 \\
治理 Skill & 引用锁定、数值守卫、权限检查和交付验证 & 规则版本、拒绝原因、回执与上游哈希 \\
\bottomrule
\end{tabularx}

\begin{minipage}[t]{0.49\textwidth}
\begin{evidencebox}{运行证据}
q001 保存 3 份假设阶段和 1 份 post-pilot 临时交互记录；选定方法 manifest 记录 Skill 来源、内容哈希与路由结果。
\end{evidencebox}
\end{minipage}\hfill
\begin{minipage}[t]{0.49\textwidth}
\begin{boundarybox}{当前边界}
受控进化已有开发/隐藏分区、影子评估和晋级代码契约；尚未发现完整真实晋级链，效果标记为\pending。
\end{boundarybox}
\end{minipage}

\subsection*{固定上下文顺序：把“记得什么”变成可审计输入}
\begin{tikzpicture}[node distance=1.5mm]
  \node[arprimary,minimum width=22mm] (c1) {角色政策};
  \node[arnode,minimum width=22mm,right=of c1] (c2) {任务字节};
  \node[arnode,minimum width=22mm,right=of c2] (c3) {允许 Skills};
  \node[arnode,minimum width=22mm,right=of c3] (c4) {锁定文献};
  \node[arnode,minimum width=22mm,right=of c4] (c5) {派生导航};
  \node[arnode,minimum width=22mm,right=of c5] (c6) {已验证 pilot};
  \node[arblock,minimum width=22mm,right=of c6] (c7) {输出 Schema};
  \draw[ardata] (c1)--(c2); \draw[ardata] (c2)--(c3); \draw[ardata] (c3)--(c4);
  \draw[ardata] (c4)--(c5); \draw[ardata] (c5)--(c6); \draw[ardata] (c6)--(c7);
\end{tikzpicture}

\small
\begin{tabularx}{\textwidth}{>{\bfseries\color{ARnavy}}p{28mm}X X}
\toprule
风险 & 代码控制 & 当前可验证对象 \\
\midrule
摘要污染事实 & 派生上下文显式标记为非证据；结论须回到原始文件 & raw-memory envelope、派生记忆字段、证据绑定 \\
Skill 越权取数 & 路由只装配容量内、角色允许的 Skill 与证据子集 & 路由理由、字符预算、Skill 内容哈希 \\
任务后残留权限 & 完成或异常时生成 cleanup/emergency receipt & q001 交互产物；批量清理成功率\pending \\
\bottomrule
\end{tabularx}
\normalsize
\clearpage

% ==================== P8 STATE MACHINE ====================
\section{十二阶段科研闭环与状态机}
\subsection{科研状态机：只有证据满足条件，状态才能跃迁}
\begin{figure}[H]
\centering
\begin{tikzpicture}[x=1cm,y=1cm]
  \node[arprimary,minimum width=2.5cm] (s0) at (-6.8,3.6) {问题已接收};
  \node[arnode,minimum width=2.5cm] (s1) at (-3.4,3.6) {证据已锁定};
  \node[arnode,minimum width=2.5cm] (s2) at (0,3.6) {候选假设已评议};
  \node[arnode,minimum width=2.5cm] (s3) at (3.4,3.6) {计划态\\plan-only};
  \node[arnode,minimum width=2.5cm] (s4) at (6.8,3.6) {可执行态\\executable};
  \node[arnode,minimum width=2.5cm] (s5) at (6.8,0.8) {执行完成};
  \node[arnode,minimum width=2.5cm] (s6) at (3.4,0.8) {观测已验证};
  \node[arnode,minimum width=2.5cm] (s7) at (0,0.8) {独立评审完成};
  \node[arnode,minimum width=2.5cm] (s8) at (-3.4,0.8) {责任授权};
  \node[arprimary,minimum width=2.5cm] (s9) at (-6.8,0.8) {正式交付/执行};

  \draw[ardata] (s0)--node[above,font=\scriptsize]{输入 Schema}(s1);
  \draw[ardata] (s1)--node[above,font=\scriptsize]{引用目录}(s2);
  \draw[ardata] (s2)--node[above,font=\scriptsize]{可证伪性}(s3);
  \draw[ardata] (s3)--node[above,font=\scriptsize]{适配器＋授权}(s4);
  \draw[ardata] (s4)--(s5); \draw[ardata] (s5)--(s6); \draw[ardata] (s6)--(s7); \draw[ardata] (s7)--(s8); \draw[ardata] (s8)--(s9);

  \node[arblock,minimum width=3.0cm] (b1) at (0,-1.8) {阻断态\\major revision / reject};
  \node[arnode,minimum width=3.2cm] (b2) at (-4.0,-1.8) {修订任务候选\\需显式确认};
  \node[arnode,minimum width=3.2cm] (b3) at (4.0,-1.8) {failure artifact\\保留错误与输入};
  \draw[arfeedback] (s7)--(b1); \draw[arfeedback] (s6)--(b3); \draw[arfeedback] (b1)--(b2);
  \draw[arfeedback] (b2) to[bend left=28] node[left,font=\scriptsize,text=ARred]{补证/修订} (s1);
  \draw[arfeedback] (b3) to[bend right=28] node[right,font=\scriptsize,text=ARred]{恢复/重试} (s4);

  \node[draw=ARred,dashed,rounded corners,fit=(s0)(s1)(s2)(s3)(s4)(s5)(s6)(s7)(b1),inner sep=3mm,label={[font=\scriptsize\sffamily,text=ARred]above:q001 已验证路径：从问题到独立评审后进入阻断态}] {};
\end{tikzpicture}
\caption{科研状态机。计划态不自动等于可执行态；科学评审和责任授权是独立跃迁条件。}
\label{fig:research-state-machine}
\end{figure}

\begin{minipage}[t]{0.49\textwidth}
\begin{evidencebox}{q001 路径｜真实运行}
问题接收 → 证据锁定 → 假设/计划 → 真实探索性 pilot → 指标验证 → 独立评审 → \texttt{major\_revision} → \blocked。没有创建正式交付报告。
\end{evidencebox}
\end{minipage}\hfill
\begin{minipage}[t]{0.49\textwidth}
\begin{boundarybox}{其余 124 题路径｜计划态}
完成问题编译与研究计划后保持 \texttt{plan-only}；由于没有兼容真实预实验适配器，不允许生成观测值或结果图。q050、q091、q119 属于该路径。
\end{boundarybox}
\end{minipage}

\subsection{状态跃迁的最小判据}
\begin{tabularx}{\textwidth}{>{\bfseries\color{ARnavy}}p{38mm}X}
\toprule
跃迁 & 必须满足 \\
\midrule
计划态 → 可执行态 & 数据授权、适配器、资源/安全审批、预注册指标和停止条件齐备 \\
执行态 → 观测已验证 & 日志、参数、环境、原始数据、metrics 与 manifest 哈希一致 \\
已评审 → 可交付 & 非阻断性科学建议＋机器校验通过＋责任主体明确授权 \\
\bottomrule
\end{tabularx}

\subsection*{状态不变量：任何重试都不能绕过}
\small
\begin{tabularx}{\textwidth}{>{\bfseries\color{ARnavy}}p{33mm}X}
\toprule
不变量 & 系统要求 \\
\midrule
证据不降级 & 派生摘要、模型回忆和未验证文本不能替代锁定文献或原始文件 \\
观测不伪造 & 计划态若无真实适配器与运行产物，只能保持 plan-only \\
失败不覆盖 & 重试创建新 checkpoint；旧 failure/review/block artifact 继续可见 \\
授权不继承 & 新仪器、数据、资源和发布动作必须重新满足对应责任门禁 \\
\bottomrule
\end{tabularx}
\normalsize

\begin{innovationbox}{状态机的实际作用}
它把“下一步该做什么”从自由对话判断改为显式条件判断：证据不足时收缩或阻断，证据充分时才允许进入执行、评审或交付。
\end{innovationbox}
\clearpage

% ==================== P9 12 STAGES ====================
\section*{7\quad 十二阶段科研闭环与状态机（续）}
\markboth{7\quad 十二阶段科研闭环与状态机（续）}{}
\subsection{十二阶段不可变 checkpoint}
\begin{figure}[H]
\centering
\begin{tikzpicture}[x=1cm,y=1cm]
  \foreach \i/\x/\label in {1/-6.6/{宽域文献检索},2/-2.2/{研究焦点收窄},3/2.2/{定向文献检索},4/6.6/{参考目录锁定}}{
    \node[arnode,minimum width=3.2cm,minimum height=1.25cm] (n\i) at (\x,3.7) {\textbf{\i}\quad \label\\[-0.5mm]{\scriptsize checkpoint＋hash}};
  }
  \foreach \i/\x/\label in {5/6.6/{Skill 路由},6/2.2/{候选假设脑暴},7/-2.2/{临时研究计划},8/-6.6/{真实 pilot / 协议}}{
    \node[arnode,minimum width=3.2cm,minimum height=1.25cm] (n\i) at (\x,1.15) {\textbf{\i}\quad \label\\[-0.5mm]{\scriptsize checkpoint＋hash}};
  }
  \foreach \i/\x/\label in {9/-6.6/{pilot 后客观复盘},10/-2.2/{证据约束修订},11/2.2/{确定性渲染},12/6.6/{独立科学评审}}{
    \node[arnode,minimum width=3.2cm,minimum height=1.25cm] (n\i) at (\x,-1.4) {\textbf{\i}\quad \label\\[-0.5mm]{\scriptsize checkpoint＋hash}};
  }
  \draw[ardata] (n1)--(n2); \draw[ardata] (n2)--(n3); \draw[ardata] (n3)--(n4);
  \draw[ardata] (n4)--(n5); \draw[ardata] (n5)--(n6); \draw[ardata] (n6)--(n7); \draw[ardata] (n7)--(n8);
  \draw[ardata] (n8)--(n9); \draw[ardata] (n9)--(n10); \draw[ardata] (n10)--(n11); \draw[ardata] (n11)--(n12);

  \node[arblock,minimum width=4.4cm] (fail) at (-3.8,-3.9) {failure artifact\\错误、输入、日志、最近可信状态};
  \node[arblock,minimum width=4.4cm] (block) at (3.8,-3.9) {review/block receipt\\问题分级、证据范围、下一轮入口};
  \draw[arfeedback] (n8)--(fail); \draw[arfeedback] (n12)--(block);
  \draw[arfeedback] (block) to[bend right=18] node[above,font=\scriptsize,text=ARred]{反馈回灌，不覆盖旧版本} (n3.south);
\end{tikzpicture}
\caption{十二阶段科研闭环。每一阶段均保存不可变 checkpoint；执行失败与评审阻断分别生成正式产物。}
\label{fig:twelve-stage-loop}
\end{figure}

\begin{innovationbox}{核心差异：过程即证据}
传统生成链常把中间推理视为可丢弃缓存；AutoResearch 把阶段输入、输出、哈希、模型回执和失败原因共同作为审计对象。下游若要改变上游结论，只能创建新版本并保留差异。
\end{innovationbox}

\begin{tabularx}{\textwidth}{>{\bfseries\color{ARnavy}}p{34mm}X X}
\toprule
反馈类型 & 例子 & 状态动作 \\
\midrule
机器可判定 & Schema 缺字段、引用越界、数字无来源、manifest 漂移 & 立即拒绝；生成定位到字段/文件的回执 \\
科学评审 & 零模型不足、方法引用不匹配、创新性检索不充分、外推过度 & major/minor 分类；阻断或创建修订入口 \\
人类责任 & 仪器、伦理、私有数据、云资源、公开发布 & 明确授权后才升级为可执行/可交付 \\
\bottomrule
\end{tabularx}

\begin{boundarybox}{当前闭环边界}
q001 已验证“真实 pilot → 独立评审 → 阻断”；本次运行没有自动跨越阻断进入“修订 → 再执行 → 再评审”。因此文档只把评审问题写成下一轮任务候选，不声称已完成多轮自动自我改进。
\end{boundarybox}

\subsection*{checkpoint 与失败回执的最小字段}
\scriptsize
\begin{tabularx}{\textwidth}{>{\bfseries\color{ARnavy}}p{27mm}X X}
\toprule
对象 & 必备字段 & 复核问题 \\
\midrule
阶段 checkpoint & stage、status、input/output hash、模型回执、时间戳 & 该阶段基于哪一版输入？是否被后续静默改写？ \\
failure artifact & 失败类型、错误信息、输入绑定、日志、最近可信状态 & 为什么失败？可以从哪里安全恢复？ \\
review/block receipt & recommendation、问题分级、证据范围、是否创建交付 & 谁判定阻断？哪些证据进入了判断？ \\
\bottomrule
\end{tabularx}
\normalsize
\clearpage

% ==================== P10 EXECUTION / GUARDS / GATES ====================
\section{真实执行、数值守卫与三层门禁}
\subsection{从计划中的变量到可复核观测}
\begin{figure}[H]
\centering
\begin{tikzpicture}[x=1cm,y=1cm]
  \node[arprimary,minimum width=3.1cm] (p) at (-6.5,2.9) {预注册计划\\变量/基线/指标};
  \node[arnode,minimum width=3.1cm] (a) at (-2.6,2.9) {领域适配器\\沙箱或受控仪器};
  \node[arnode,minimum width=3.1cm] (r) at (1.3,2.9) {原始产物\\数据/日志/环境};
  \node[arnode,minimum width=3.1cm] (m) at (5.2,2.9) {metrics＋manifest\\文件 SHA-256};
  \draw[ardata] (p)--(a); \draw[ardata] (a)--(r); \draw[ardata] (r)--(m);

  \node[arnode,minimum width=3.1cm] (ph) at (-6.5,0.3) {报告占位符\texttt{[[pilot:path]]}};
  \node[arnode,minimum width=3.1cm] (ng) at (-2.6,0.3) {数值守卫\\仅取已验证字段};
  \node[arnode,minimum width=3.1cm] (rl) at (1.3,0.3) {引用目录锁\\精确投影};
  \node[arblock,minimum width=3.1cm] (out) at (5.2,0.3) {证据约束修订\\无来源即拒绝};
  \draw[ardata] (ph)--(ng); \draw[ardata] (m)--(ng); \draw[ardata] (ng)--(rl); \draw[ardata] (rl)--(out);

  \node[arnode,minimum width=3.6cm,minimum height=1.3cm] (g1) at (-4.7,-2.6) {\textbf{机器门禁}\\Schema / 数字 / 引用 / manifest};
  \node[arnode,minimum width=3.6cm,minimum height=1.3cm] (g2) at (0,-2.6) {\textbf{科学门禁}\\独立、非重写评审};
  \node[arnode,minimum width=3.6cm,minimum height=1.3cm] (g3) at (4.7,-2.6) {\textbf{人类门禁}\\责任、伦理、资源与发布授权};
  \draw[ardata] (g1)--(g2); \draw[ardata] (g2)--(g3);
  \draw[arfeedback] (out)--(g1);
  \node[arblock,minimum width=3.0cm] (stop) at (0,-4.55) {任一阻断性失败\\不得自动交付};
  \draw[arfeedback] (g1)--(stop); \draw[arfeedback] (g2)--(stop); \draw[arfeedback] (g3)--(stop);
\end{tikzpicture}
\caption{真实执行、数值/引用守卫与三层门禁。结构通过不代表科学通过，科学通过也不代表责任授权。}
\end{figure}

\begin{minipage}[t]{0.49\textwidth}
\begin{evidencebox}{已发生：q001}
真实计算产生原始 CSV、四类零模型 CSV、日志、参数、环境、metrics 与 manifest；修订计划中的关键数字从验证文件绑定。独立评审为 \texttt{major\_revision}，随后生成阻断回执，\texttt{delivery\_report\_created=false}。
\end{evidencebox}
\end{minipage}\hfill
\begin{minipage}[t]{0.49\textwidth}
\begin{boundarybox}{尚未量化}
批量数字扰动检出率、引用越界率、数值守卫误报率、人类签署样本与三层门禁时延均为\pending。存在代码控制点不等于已经获得统计性能结论。
\end{boundarybox}
\end{minipage}

\subsection{三层门禁的职责边界}
\small
\begin{tabularx}{\textwidth}{>{\bfseries\color{ARnavy}}p{23mm}X X p{33mm}}
\toprule
门禁 & 判断问题 & 典型输出 & 当前证据 \\
\midrule
机器 & 文件和字段是否完整、数字/引用是否可追溯？ & pass / field-level reject / manifest mismatch & 代码＋q001 运行 \\
科学 & 方法、零模型、统计解释和创新性是否成立？ & accept / minor / major / reject & q001 major revision \\
人类 & 是否授权真实仪器、伦理数据、资源与正式发布？ & approve / hold / deny & 流程接口；签署样本待补 \\
\bottomrule
\end{tabularx}
\normalsize
\clearpage

% ==================== P11 SCIENCE 125 ====================
\section{Science 125 问成果总览}
\subsection{规模化交付的真实状态，而非“125 次实验”}
\begin{figure}[H]
\centering
\begin{tikzpicture}[x=1cm,y=1cm]
  \node[arprimary,minimum width=4.1cm,minimum height=1.8cm] (c1) at (-5.4,3.2) {\fontsize{22}{24}\selectfont\bfseries 125/125\\[-1mm]{\small 结构化研究计划}};
  \node[arnode,minimum width=4.1cm,minimum height=1.8cm] (c2) at (0,3.2) {\fontsize{22}{24}\selectfont\bfseries 125\\[-1mm]{\small Markdown 交付}};
  \node[arnode,minimum width=4.1cm,minimum height=1.8cm] (c3) at (5.4,3.2) {\fontsize{22}{24}\selectfont\bfseries 125\\[-1mm]{\small PDF 交付}};

  \node[arnode,minimum width=9.9cm,minimum height=1.25cm,anchor=west] (plan) at (-7.45,0.8) {\textbf{计划态 124 题}\quad q002--q125：完整协议，明确“尚未执行预实验”};
  \node[arprimary,minimum width=4.8cm,minimum height=1.25cm,anchor=west] (pilot) at (2.65,0.8) {\textbf{真实 pilot 1 题}\quad q001：探索性计算};
  \draw[ardata] (c1.south)--(plan.north); \draw[ardata] (c1.south)--(pilot.north);

  \node[arnode,minimum width=3.35cm,minimum height=1.55cm] (q1) at (-6.0,-1.7) {q001\\计算数论\\{\scriptsize 已执行 pilot}};
  \node[arnode,minimum width=3.35cm,minimum height=1.55cm] (q50) at (-2.0,-1.7) {q050\\宇宙学\\{\scriptsize 计划态}};
  \node[arnode,minimum width=3.35cm,minimum height=1.55cm] (q91) at (2.0,-1.7) {q091\\超导信息器件\\{\scriptsize 计划态}};
  \node[arnode,minimum width=3.35cm,minimum height=1.55cm] (q119) at (6.0,-1.7) {q119\\认知科学\\{\scriptsize 计划态}};
  \node[font=\scriptsize\sffamily,text=ARgray] at (0,-3.0) {四例仅说明跨领域任务化能力；125 题完整领域计数尚未形成可审计统计，标记为【待补统计】。};
\end{tikzpicture}
\caption{Science 125 问成果仪表板。覆盖、交付格式和真实执行状态分开统计。}
\label{fig:science125-dashboard}
\end{figure}

\begin{tabularx}{\textwidth}{>{\bfseries\color{ARnavy}}p{16mm}p{34mm}X p{31mm}}
\toprule
编号 & 领域问题 & 计划中的关键对照与判据 & 状态 \\
\midrule
q001 & 素数间隙序贯结构 & 五个固定区间；四类零模型；$m=5$ 并列感知排列熵 & \runproof：探索性 pilot \\
q050 & Hubble 张力与 FLRW 破缺 & 时间延迟、BAO、SN Ia；$>3\sigma$、$\ln B>5$、夹角 $<30^\circ$ & 计划态；未执行 \\
q091 & RQL 高频上限 & 1--60 GHz 扫频；结参数扰动、磁通干预、BER/噪声带宽 & 计划态；未执行 \\
q119 & 工作记忆槽位可塑性 & 四周纵向、双任务/干扰；$d$、$p$、$\Delta$BIC & 计划态；未执行 \\
\bottomrule
\end{tabularx}

\begin{innovationbox}{跨领域规模化任务编译}
每题将开放问题压缩为可检查对象：Problem、Rationale、Technical Details、Datasets、Methods、Experiments、Results 与 References，并补充基线、指标、终止条件、资源和边界。125/125 的当前实测是\textbf{必备章节覆盖}；JSON Schema 级批量完整率仍为\pending。
\end{innovationbox}

\begin{boundarybox}{必须与成果包共同阅读的边界}
“pilot 方案”表示可执行研究协议，不表示已产生观测。只有 q001 当前具有可核验的原始数据、零模型数据、日志、指标和 manifest。
\end{boundarybox}

\subsection*{每题交付对象的最小可执行语义}
\scriptsize
\begin{tabularx}{\textwidth}{>{\bfseries\color{ARnavy}}p{31mm}X X}
\toprule
对象组 & 必备内容 & 允许的状态 \\
\midrule
科学目标 & 分析单位、主假设、竞争解释、最强反证 & draft / evidence-locked \\
执行计划 & 数据源、变量、基线、方法、指标、终止条件、资源 & plan-only / executable \\
观测与证据 & 原始文件、参数、环境、日志、metrics、manifest & absent / executed / verified \\
决策与交付 & 评审、版本差异、风险、授权与阻断原因 & revise / blocked / deliverable \\
\bottomrule
\end{tabularx}
\normalsize
\clearpage

% ==================== P12 Q001 ====================
\section{q001 真实 pilot：可信闭环压力测试}
\subsection{从开放问题到“应当收缩结论”的真实反馈}
\begin{figure}[H]
\centering
\begin{tikzpicture}[node distance=1mm]
  \node[arprimary,minimum width=20mm] (q) {问题\\素数结构};
  \node[arnode,minimum width=21mm,right=of q] (h) {有限假设\\序贯复杂度};
  \node[arnode,minimum width=28mm,right=of h] (n) {四类零模型\\局部/全局/残基/wheel};
  \node[arnode,minimum width=22mm,right=of n] (x) {真实计算\\5 固定区间};
  \node[arnode,minimum width=22mm,right=of x] (m) {指标锁定\\metrics＋hash};
  \node[arnode,minimum width=23mm,right=of m] (r) {独立评审\\major revision};
  \node[arblock,minimum width=22mm,right=of r] (b) {阻断\\不交付};
  \draw[ardata] (q)--(h); \draw[ardata] (h)--(n); \draw[ardata] (n)--(x); \draw[ardata] (x)--(m); \draw[ardata] (m)--(r); \draw[ardata] (r)--(b);
  \node[arnode,minimum width=6.8cm,below=7mm of m] (next) {下一轮任务候选：更正方法引用｜补相邻工作｜统一 mod-30 / wheel-210 表述｜保持数值精度};
  \draw[arfeedback] (b.south) to[bend left=12] (next.east);
  \draw[arfeedback] (next.west) to[bend left=14] node[below,font=\scriptsize,text=ARred]{需人确认后重入证据阶段} (h.south);
\end{tikzpicture}
\caption{q001 案例链路。系统保存弱化结论和负面评审，而非把探索性结果包装为突破。}
\label{fig:q001-chain}
\end{figure}

\begin{minipage}[t]{0.58\textwidth}
\small
\begin{tabularx}{\textwidth}{>{\bfseries\color{ARnavy}}X r r r}
\toprule
零模型 & 零模型均值 & 观测－零模型 & Holm $p$ \\
\midrule
局部分块置换 & 0.953933 & $-0.024580$ & 0.020 \\
全局置换 & 0.954408 & $-0.025055$ & 0.020 \\
残基路径条件置换 & 0.930507 & $-0.001153$ & 0.020 \\
wheel-210 约束 & 0.931604 & $-0.002251$ & 0.020 \\
\bottomrule
\end{tabularx}
\normalsize
\begin{evidencebox}{结果如何改变研究计划}
观测均值为 0.929353。相对全局置换的差异在算术条件零模型下缩小约一个数量级，提示大部分表观结构可由已知模约束解释。剩余差异只构成扩大尺度和改进零模型的线索。
\end{evidencebox}
\end{minipage}\hfill
\begin{minipage}[t]{0.39\textwidth}
\begin{reviewbox}{独立评审：major revision}
\textbf{Major 1}\quad 方法引用与频率学派、并列感知实现不完全匹配。\par
\textbf{Major 2}\quad 创新性三角检索与相邻工作比较不足。\par
\textbf{Minor}\quad mod-30 / wheel-210 叙述不一致；小数精度需统一。\par
\medskip\textbf{系统动作：}\blocked；未生成自动交付报告。
\end{reviewbox}
\end{minipage}

\begin{boundarybox}{科学解释边界}
五个区间是预先固定的计算 benchmark，不是从数轴总体随机抽取的样本；经验 $p$ 与区间重采样范围仅描述本次计算。该 pilot \textbf{没有证明或否定任何素数猜想}，也不构成正式数学证明。
\end{boundarybox}

\begin{evidencebox}{可复核证据链}
5 份原始区间 CSV＋5 份零模型 CSV＋环境/参数＋stdout/stderr＋metrics＋manifest＋源计划/代码快照；独立评审记录 18 个 verified files。完整路径与文件索引见附录。
\end{evidencebox}

\subsection*{下一轮候选任务与再进入条件}
\small
\begin{tabularx}{\textwidth}{>{\bfseries\color{ARnavy}}p{31mm}X X}
\toprule
评审问题 & 下一轮任务候选 & 重新进入评审前的检查 \\
\midrule
方法引用不匹配 & 补入与实际频率学派、并列感知排列熵一致的方法来源 & 方法描述、代码实现与引用三者逐项对应 \\
创新性证据不足 & 扩展相邻工作检索并给出实质差异，不只列题名 & 锁定目录更新；主张均有直接支持或收缩 \\
模数叙述不一致 & 统一残基路径与 wheel-210 的实际实现及文字边界 & 参数、代码、图表和正文使用同一术语 \\
数值精度漂移 & 从 metrics 重新绑定表格与正文，不手工改写 & 数字守卫和 manifest 复核通过 \\
\bottomrule
\end{tabularx}
\normalsize

\begin{boundarybox}{当前状态}
以上为由评审问题导出的\textbf{候选修订任务}；本次运行在阻断态停止，尚未自动执行这些任务，也未完成再评审。
\end{boundarybox}
\clearpage

% ==================== P13 EVALUATION ====================
\section{系统评估、消融与竞品对比}
\subsection{把“已有实测”与“计划评估”分开}
\small
\begin{tabularx}{\textwidth}{>{\bfseries\color{ARnavy}}p{37mm}p{35mm}X}
\toprule
指标/事实 & 当前结果 & 解释边界 \\
\midrule
计划与交付覆盖 & 125/125 计划；125 MD＋125 PDF & 数量不代表实验执行 \\
必备章节完整 & 125/125 & 是文档章节检查，不等于 JSON Schema 通过率 \\
真实执行覆盖 & q001：1；其余计划态：124 & 当前不是 125 项实验 \\
q001 证据文件 & 18 个 verified files & 仅对本次独立评审范围 \\
q001 引用集合 & \texttt{verified\_exact\_union} & 仅对该锁定目录与产物 \\
科学门禁事件 & 1 次 major revision，自动交付被阻断 & 单样本事实，不外推“阻断率” \\
\bottomrule
\end{tabularx}
\normalsize

\subsection{待补实验矩阵}
\scriptsize
\begin{tabularx}{\textwidth}{>{\bfseries\color{ARnavy}}p{25mm}p{34mm}p{29mm}p{32mm}X}
\toprule
实验组 & 目的 / 对照 & 自变量 & 指标 / 样本 & 回答的问题 \\
\midrule
Schema＋manifest & 验证对象与多格式一致；对照＝无校验 & 校验器开/关、受控篡改 & 通过率、漂移检出；125 套 & 规模化输出是否真正遵守契约？ \\
数字/引用对抗 & 检验守卫；对照＝自由生成 & 错误类型/幅度、守卫开关 & TPR/FPR、定位率；≥300 扰动 & 能否拒绝无来源数字与越界引用？ \\
计划/观测隔离 & 防止未执行结果化；对照＝无状态约束 & 契约/校验开关 & 混淆率；124 题＋对抗提示 & 能否稳定保持 plan-only？ \\
checkpoint 故障注入 & 验证恢复；对照＝全流程重跑 & 12 个故障阶段、恢复策略 & 恢复率、重复量、一致性；每阶段≥5 & 长流程是否可从可信状态续跑？ \\
反馈质量曲线 & 验证迭代增益；对照＝单轮 & 0/1/2/3 轮、反馈来源 & 专家盲分、缺陷、工时；20 题 & 反馈何时提升、何时饱和？ \\
模型/Skills 消融 & 分离系统贡献；对照＝直接 Qwen & RAG、Agent、Skill、门禁组合 & 质量/溯源/成本；20--30 题 & 增益来自模型还是治理机制？ \\
Skill 受控进化 & 验证影子晋级；对照＝直接写回/不写回 & 隐藏评估、晋级策略 & 提升/回归/回滚；≥20 候选 & 能否学习而不污染安全与证据？ \\
成本与人因 & 建立可落地预算；对照＝直接 LLM/RAG & 模型、上下文、并发 & Token、费用、P50/P95、修改次数；≥30 题 & 质量约束下的实际代价是什么？ \\
\bottomrule
\end{tabularx}
\normalsize

\begin{innovationbox}{最小可执行消融设计}
在 20 个跨学科问题上比较：A 直接 Qwen；B Qwen＋锁定检索；C B＋临时智能体/Skills；D 完整 AutoResearch（状态契约＋守卫＋门禁）。采用同一输入、预算与专家盲评量表，同时报告质量、证据错误、阻断表现、时间和成本。当前结果统一标记为\pending。
\end{innovationbox}

\begin{boundarybox}{评估纪律}
本页没有用代码存在性代替效果数据，也没有生成虚假百分比、成本或质量提升。完整实验目的、对照、自变量、指标、样本和研究问题见配套 Markdown 审计稿。
\end{boundarybox}

\subsection*{赛前补证优先级}
\begin{minipage}[t]{0.24\textwidth}\begin{flatbox}\textbf{P0｜125 题契约复核}\par 全量 Schema＋manifest；先回答“规模化是否真实合规”。\end{flatbox}\end{minipage}\hfill
\begin{minipage}[t]{0.24\textwidth}\begin{flatbox}\textbf{P0｜数字/引用对抗}\par 受控注入 ≥300 个扰动；报告 TPR、FPR 和定位率。\end{flatbox}\end{minipage}\hfill
\begin{minipage}[t]{0.24\textwidth}\begin{flatbox}\textbf{P1｜故障恢复}\par 12 阶段各 ≥5 次注入；报告恢复率与状态一致性。\end{flatbox}\end{minipage}\hfill
\begin{minipage}[t]{0.24\textwidth}\begin{flatbox}\textbf{P1｜20 题消融}\par 同输入、同预算、专家盲评；分离模型与治理贡献。\end{flatbox}\end{minipage}

\begin{innovationbox}{完成判据}
只有在实验脚本、原始结果、统计汇总、失败样本与配置清单同时归档后，待补项才可从“实验设计”升级为“实测结论”。
\end{innovationbox}
\clearpage

% ==================== P14 COMPETITORS / ENGINEERING ====================
\section*{11\quad 系统评估、消融与竞品对比（续）}
\markboth{11\quad 系统评估、消融与竞品对比（续）}{}
\subsection{基于论文与官方仓库的克制对照}
\scriptsize
\begin{tabularx}{\textwidth}{>{\bfseries\color{ARnavy}}p{27mm}p{42mm}X}
\toprule
方案 & 公开材料明确的主要机制 & 与本方案相关的审计边界 \\
\midrule
普通大模型直接生成 & 基于输入上下文生成文本 & 不构成统一产品，证据状态、checkpoint 和科研门禁取决于具体实现 \\
RAG 科研助手\cite{rag} & 检索器＋生成器，将外部知识引入生成\cite{rag} & 原始 RAG 论文未定义计划/观测状态、数值守卫或科学责任门禁 \\
通用多智能体（AutoGen）\cite{autogen} & 消息传递、事件驱动 Agent、人类与工具协作 & 官方框架面向通用编排；本文所列科研治理机制在公开材料中未明确为默认能力 \\
AI Scientist\cite{aiscientist} & 在 ML 模板中完成想法、实验、论文写作与自动评审 & 官方仓库明确三类核心模板；本文的证据对象/数值守卫/三层门禁未在公开材料中明确说明 \\
AI Scientist-v2\cite{aiscientist2} & 移除人工代码模板，使用实验经理与渐进式 Agent 树搜索 & 论文报告端到端 ML 研究；本文治理维度若未披露，不推断为“没有” \\
AutoResearch & 科研状态契约、临时 Agent/Skills、不可变闭环、守卫与门禁、125 问任务编译 & q001 有运行证据；部分系统指标与 Skill 晋级仍待实验 \\
\bottomrule
\end{tabularx}
\normalsize

\begin{boundarybox}{比较结论}
该表比较的是\textbf{公开披露的系统重心与审计机制}，不是同预算性能排行榜。没有共同任务实验时，不宣称 AutoResearch 全面优于既有系统；“公开材料中未明确说明”也不等于该系统绝对不具备相应能力。
\end{boundarybox}

\section{工程实现、部署与复现}
\begin{minipage}[t]{0.32\textwidth}
\begin{evidencebox}{模型与工具接口}
Qwen 为当前推理核心；兼容接口从配置读取端点、密钥和模型名。工具调用、领域适配器与模型层解耦。\par
\textbf{状态：}\codeproof；跨供应商基准\pending。
\end{evidencebox}
\end{minipage}\hfill
\begin{minipage}[t]{0.32\textwidth}
\begin{evidencebox}{知识与版本底座}
Obsidian 兼容 vault 保存文献、实验、问题、失败、Skill 与策略卡；原始证据与派生导航分层，Git 提供版本历史。\par
\textbf{状态：}\codeproof＋运行文件。
\end{evidencebox}
\end{minipage}\hfill
\begin{minipage}[t]{0.32\textwidth}
\begin{evidencebox}{确定性交付}
同一结构化计划渲染 Markdown、TeX、PDF 与 manifest；源哈希、证据绑定哈希和产物清单用于复核。\par
\textbf{状态：}\runproof＋\codeproof。
\end{evidencebox}
\end{minipage}

\subsection{最小复现顺序}
\begin{tikzpicture}[node distance=3mm]
  \node[arprimary,minimum width=27mm] (r1) {1 锁定依赖\\与模型配置};
  \node[arnode,minimum width=27mm,right=of r1] (r2) {2 编译题目\\锁定文献目录};
  \node[arnode,minimum width=27mm,right=of r2] (r3) {3 路由 Skills\\生成临时计划};
  \node[arnode,minimum width=27mm,right=of r3] (r4) {4 有适配器则\\执行真实 pilot};
  \node[arnode,minimum width=27mm,right=of r4] (r5) {5 验证 manifest\\绑定数字/引用};
  \node[arblock,minimum width=27mm,right=of r5] (r6) {6 独立评审\\交付或阻断};
  \draw[ardata] (r1)--(r2); \draw[ardata] (r2)--(r3); \draw[ardata] (r3)--(r4); \draw[ardata] (r4)--(r5); \draw[ardata] (r5)--(r6);
\end{tikzpicture}

\begin{boundarybox}{部署边界}
默认沙箱；不在未授权情况下操作真实仪器、访问私有数据、租用云 GPU、提交论文或公开发布。凭据由环境注入，不写入运行目录、知识库和 Git 历史。
\end{boundarybox}

\subsection*{复现包的最小清单}
\small
\begin{tabularx}{\textwidth}{>{\bfseries\color{ARnavy}}p{29mm}X X}
\toprule
复现层 & 必备对象 & 复核目标 \\
\midrule
环境 & 依赖锁、模型配置名、适配器版本、随机种子、硬件/系统快照 & 能否解释版本与环境差异 \\
输入 & 题目快照、锁定文献目录、数据授权、Skill manifest、输入哈希 & 能否重建任务边界 \\
执行 & 原始数据、参数、stdout/stderr、metrics、文件清单与 SHA-256 & 能否复算并发现篡改 \\
决策 & 修订对象、独立评审、阻断/授权回执、parent hash 与版本差异 & 能否解释为何交付或停止 \\
\bottomrule
\end{tabularx}
\normalsize

\begin{evidencebox}{当前可复现样本}
q001 已具备输入、运行、指标、manifest、独立评审和阻断回执；跨机器一键复现率与环境差异统计仍为\pending。
\end{evidencebox}
\clearpage

% ==================== P15 SAFETY / VALUE / ROADMAP ====================
\section{安全边界、应用价值与下一阶段计划}
\subsection{安全不是附加声明，而是状态跃迁条件}
\begin{minipage}[t]{0.32\textwidth}\begin{flatbox}\textbf{最小权限}\par Agent 只读必要 Skill 与证据子集；工具按白名单开放。\end{flatbox}\end{minipage}\hfill
\begin{minipage}[t]{0.32\textwidth}\begin{flatbox}\textbf{证据不可覆盖}\par 原始字节、日志、失败和旧版本保留；摘要不得升级为证据。\end{flatbox}\end{minipage}\hfill
\begin{minipage}[t]{0.32\textwidth}\begin{flatbox}\textbf{危险动作审批}\par 仪器、伦理数据、云资源和正式发布需要人类授权。\end{flatbox}\end{minipage}

\begin{minipage}[t]{0.32\textwidth}\begin{flatbox}\textbf{密钥隔离}\par 模型凭据只从环境读取，不进入产物、vault 与仓库。\end{flatbox}\end{minipage}\hfill
\begin{minipage}[t]{0.32\textwidth}\begin{flatbox}\textbf{受控进化}\par 候选 Skill 先影子评估；安全、审批与发布规则不可自改。\end{flatbox}\end{minipage}\hfill
\begin{minipage}[t]{0.32\textwidth}\begin{flatbox}\textbf{责任可追溯}\par 自动评分只筛选；最终科学判断和高风险行为归责任主体。\end{flatbox}\end{minipage}

\subsection{应用价值}
\begin{tabularx}{\textwidth}{>{\bfseries\color{ARnavy}}p{31mm}X X}
\toprule
使用场景 & 可交付价值 & 不替代的责任 \\
\midrule
前沿问题预研 & 把宽问题快速编译为带证据、基线、指标、停止条件的方案集合 & 领域专家的科学判断与研究优先级 \\
计算/仿真研究 & 连接适配器后形成原始产物—指标—评审的可复核链 & 模型有效性、资源授权与正式结论 \\
实验协同与项目治理 & 保存进度、失败、版本差异和下一轮任务，减少信息断裂 & 伦理审批、仪器安全与发布责任 \\
\bottomrule
\end{tabularx}

\subsection{下一阶段路线图}
\begin{tikzpicture}[x=1cm,y=1cm]
  \node[arprimary,minimum width=5.0cm,minimum height=2.0cm] (p1) at (-5.5,0) {\textbf{阶段 I｜赛前补证}\quad 0--4 周\\[-0.5mm]\small 125 题 Schema/manifest 批量复核\\数字/引用对抗集；20 题最小消融};
  \node[arnode,minimum width=5.0cm,minimum height=2.0cm] (p2) at (0,0) {\textbf{阶段 II｜闭合反馈}\quad 1--3 月\\[-0.5mm]\small q001 修订—再运行—再评审\\checkpoint 故障注入；人类门禁样本};
  \node[arnode,minimum width=5.0cm,minimum height=2.0cm] (p3) at (5.5,0) {\textbf{阶段 III｜领域扩展}\quad 3--12 月\\[-0.5mm]\small 为高价值题接入数据/仿真/仪器\\跨实验室复现与成本评估};
  \draw[ardata] (p1)--(p2); \draw[ardata] (p2)--(p3);
\end{tikzpicture}

\section*{结论}
\addcontentsline{toc}{section}{结论}
AutoResearch 已完成的不是“125 次自动科学发现”，而是\textbf{125 个可继续推进的结构化研究任务}，以及一个能够暴露自身缺陷的 q001 真实闭环样本。其技术实质是让 Qwen 在证据对象、最小权限 Skills、不可变 checkpoint、数字/引用守卫和分层门禁中承担规模化认知劳动。q001 的 major revision 与交付阻断表明：系统的目标不是保证每次生成成功，而是保证科研状态不越过证据边界。

\begin{reviewbox}{给评委的最终判断依据}
若只看文字生成质量，本项目容易与 RAG 或多智能体系统混淆；若检查对象哈希、计划/观测分离、真实 pilot 文件、独立评审和阻断回执，可以看到其核心差异：\textbf{自动化必须经过可审计、可失败的科研门禁。}
\end{reviewbox}

\subsection*{阶段验收标准}
\small
\begin{tabularx}{\textwidth}{>{\bfseries\color{ARnavy}}p{28mm}X X}
\toprule
阶段 & 进入下一阶段的必要证据 & 不满足时的动作 \\
\midrule
赛前补证完成 & 批量校验与对抗实验均有脚本、原始结果、统计摘要和 manifest & 保留【待补实验】，不得把设计写成指标 \\
反馈闭合完成 & q001 修订依据、再运行文件、第二次独立评审和版本差异齐备 & 维持 blocked / revise，不创建正式交付 \\
领域扩展完成 & 目标题获得数据/伦理/仪器授权并通过适配器验证与专家复核 & 维持 plan-only，优先输出待执行协议 \\
\bottomrule
\end{tabularx}
\normalsize

\begin{boundarybox}{长期边界}
系统可以扩大科研认知劳动的吞吐量，但不能把自动化速度当作科学真值；安全策略、责任门禁、证据定义和发布规则始终保留人类监督与版本回滚。
\end{boundarybox}
\clearpage

% ==================== P16 APPENDIX / REFERENCES ====================
\section*{附录：证据索引、复核路径与参考文献}
\addcontentsline{toc}{section}{附录：证据索引、复核路径与参考文献}
\markboth{附录与参考文献}{}

\subsection*{A. 核心代码与机制映射}
\scriptsize
\begin{tabularx}{\textwidth}{>{\raggedright\arraybackslash}p{72mm}X}
\toprule
相对仓库路径 & 审计职责 \\
\midrule
\path{src/autoresearch/competition/contest_direct_plan.py} & 科研计划与产物 Schema、输入/产物哈希 \\
\path{src/autoresearch/competition/temporary_qwen_pool.py} & 临时 Qwen Agent、角色/Skill/证据预算、清理回执 \\
\path{src/autoresearch/competition/contest_direct_skill_router.py} & Skill 证据路由、来源、理由与内容哈希 \\
\path{src/autoresearch/competition/contest_direction_skill_evolution.py} & 开发/隐藏分区、影子评估、晋级与拒绝契约 \\
\path{src/autoresearch/knowledge/raw_memory.py} & 原始证据不可变信封与哈希验证 \\
\path{src/autoresearch/competition/contest_direction_memory.py} & 派生记忆非证据声明与可重建导航 \\
\path{src/autoresearch/competition/contest_direction_research_loop_cli.py} & 十二阶段 checkpoint、阻断性评审状态 \\
\path{src/autoresearch/competition/contest_direct_plan_revision.py} & pilot 占位绑定、已验证文件与数值守卫 \\
\path{src/autoresearch/competition/contest_reference_policy.py} & 锁定引用目录与未知引用 fail-closed \\
\path{src/autoresearch/competition/science125_batch.py} & 125 题计划态批处理与“未执行”披露 \\
\path{src/autoresearch/competition/contest_direct_plan_render.py} & Markdown/TeX/PDF/manifest 同源渲染与清单复核 \\
\bottomrule
\end{tabularx}

\subsection*{B. 运行与成果路径}
\begin{tabularx}{\textwidth}{>{\raggedright\arraybackslash}p{72mm}X}
\toprule
路径（本地审计环境） & 内容 \\
\midrule
\path{C:/Users/Z/Desktop/}\allowbreak 揭榜挂帅提交材料\path{/delivery-125-plan} & q001--q125：125 份 Markdown＋125 份 PDF \\
\path{E:/ai-researcher-loop/ai-researcher-loop-main/}\\\path{runs/contest-delivery/live-verify-q001-v2-20260816} & q001 完整运行、checkpoint、评审与阻断回执 \\
\quad\path{preexperiment/prime-gap-information-theory-v1} & 5 份 raw CSV、5 份 null CSV、环境、参数、日志、metrics、manifest \\
\quad\path{independent-scientific-review/}\\\path{system-plan-scientific-review.json} & 18 个 verified files、引用精确并集、major revision \\
\quad\path{scientific-review-blocked-receipt.json} & \texttt{blocked\_by\_independent\_scientific\_review} \\
\bottomrule
\end{tabularx}

\subsection*{C. q001 数值与样本边界}
观测均值 0.929353；局部、全局、残基路径、wheel-210 四类零模型均值依次为 0.953933、0.954408、0.930507、0.931604；差值依次为 $-0.024580$、$-0.025055$、$-0.001153$、$-0.002251$。每类零模型 199 次重采样；五个固定 benchmark 区间的间隙数为 70,434、64,335、61,937、59,335、56,359。完整区间级统计以 \texttt{metrics.json} 为准。

\subsection*{参考文献}
\begin{thebibliography}{9}\setlength{\itemsep}{0.4pt}\setlength{\parskip}{0pt}
\bibitem{science125} Kennedy, D.; Norman, C. What Don't We Know? \textit{Science} 309, 75 (2005).
\bibitem{qwen3} Yang, A. et al. Qwen3 Technical Report. arXiv:2505.09388 (2025); official repository: \url{https://github.com/QwenLM/Qwen3}.
\bibitem{rag} Lewis, P. et al. Retrieval-Augmented Generation for Knowledge-Intensive NLP Tasks. \textit{NeurIPS} 33 (2020), arXiv:2005.11401.
\bibitem{autogen} Wu, Q. et al. AutoGen: Enabling Next-Gen LLM Applications via Multi-Agent Conversation. arXiv:2308.08155 (2023); official repository: \url{https://github.com/microsoft/autogen}.
\bibitem{aiscientist} Lu, C. et al. The AI Scientist: Towards Fully Automated Open-Ended Scientific Discovery. arXiv:2408.06292 (2024); official repository: \url{https://github.com/SakanaAI/AI-Scientist}.
\bibitem{aiscientist2} Yamada, Y. et al. The AI Scientist-v2: Workshop-Level Automated Scientific Discovery via Agentic Tree Search. arXiv:2504.08066 (2025); official repository: \url{https://github.com/SakanaAI/AI-Scientist-v2}.
\bibitem{bandt} Bandt, C.; Pompe, B. Permutation Entropy: A Natural Complexity Measure for Time Series. \textit{Phys. Rev. Lett.} 88, 174102 (2002).
\bibitem{primebias} Lemke Oliver, R. J.; Soundararajan, K. Unexpected Biases in the Distribution of Consecutive Primes. \textit{PNAS} 113, E4446--E4454 (2016).
\bibitem{provo} W3C. PROV-O: The PROV Ontology. W3C Recommendation (2013), \url{https://www.w3.org/TR/prov-o/}.
\end{thebibliography}

\vfill
\begin{boundarybox}{材料真实性声明}
本文依据榜题、源代码、运行 checkpoint、q001 真实 pilot 与 125 题既有交付包重构；未生成新的科学观测，未把计划态写成已执行结果，未删除负面评审，也未把 q001 表述为数学证明。外部系统比较仅使用论文、官方仓库或正式文档。
\end{boundarybox}

\end{document}
